\section{Introduction}\label{sec:introduction}

How many evaluations of a potential and its gradient are needed
to sample its Gibbs distribution? Even when the potential is
uniformly strongly convex and smooth, familiar discretizations
incur an error that grows with dimension. The question is
particularly sharp when conditioning and accuracy are held fixed:
is a positive power of the dimension intrinsically necessary,
or is it a feature of the sampling algorithms used?

We answer this question in the exact first-order oracle model
with unrestricted real computation between queries. The optimal
dimension exponent is zero. The construction uses approximation
orders that can depend on the desired exponent and is aimed at
oracle complexity rather than practical running time.

\subsection{Model and result}

For \(d\ge2\), let \(\mathcal V_d\) consist of the functions
\(V\in C^2(\R^d)\) satisfying
\begin{equation}\label{eq:original-class}
 V(0)=0,\qquad \nabla V(0)=0,\qquad
 I_d\preceq\nabla^2V(x)\preceq2I_d
                         \quad\text{for every }x\in\R^d.
\end{equation}
The target law is
\[
 \pi_V(dx)=Z_V^{-1}e^{-V(x)}\,dx,\qquad
 Z_V=\int_{\R^d}e^{-V(x)}\,dx.
\]
An exact first-order query at \(x\) returns
\((V(x),\nabla V(x))\). Queries may be adaptive.
Random real variables may be generated, and computation
between queries is unrestricted. An algorithm may depend
on \(d\), but its only information about \(V\) is through
the oracle responses. All algorithms are measurable.

Write \(\norm{\mu-\eta}_{\TV}=\sup_B|\mu(B)-\eta(B)|\).
Let \(Q(d)\) be the smallest integer \(q\) for which
an algorithm makes at most \(q\) queries on \emph{every}
execution and returns \(X\) satisfying
\(\norm{\mathcal L(X)-\pi_V}_{\TV}\le1/10\) for
every \(V\in\mathcal V_d\). Thus the budget is not
merely an expectation. Define the optimal dimension exponent
by
\[
 \gamma_*=\inf\{\gamma\ge0:Q(d)\le d^{\gamma+o(1)}\}
          =\limsup_{d\to\infty}\frac{\log Q(d)}{\log d}.
\]

\begin{theorem}\label{thm:main}
In the oracle model above, for every \(\varepsilon>0\)
there is \(C_\varepsilon<\infty\) such that
\[
                  Q(d)\le C_\varepsilon d^\varepsilon
                  \qquad(d\ge2).
\]
There are also absolute constants \(c>0\) and \(d_0\)
such that
\[
                  Q(d)\ge c\log d\qquad(d\ge d_0).
\]
The lower bound applies to arbitrary randomized adaptive
first-order algorithms. In particular, \(\gamma_*=0\).
\end{theorem}

The upper assertion is proved in Theorem~\ref{thm:upper}
and the lower assertion in Theorem~\ref{thm:lower}.
Together they determine the exponent but not the finer
growth rate of \(Q(d)\). In particular, the upper bound
does not assert a polylogarithmic query count.
The constants and the algorithmic approximation orders
may depend on \(\varepsilon\). No arithmetic or
bit-complexity bound is imposed or deduced.

\subsection{Historical context}

Sampling by discretizing Langevin or Hamiltonian dynamics turns
functional inequalities and smoothness into algorithmic bounds, but it
also introduces a dimension-dependent discretization problem.
Nonasymptotic analyses of the unadjusted Langevin algorithm were developed
by \citet{dalalyan2017langevin} and \citet{durmus2017langevin};
\citet{cheng2018underdamped} analyzed the underdamped method under strong
convexity and smoothness. Randomized midpoint discretization further
improved the quadratic-Wasserstein guarantees
\citep{shen2019randomized}. Wasserstein and total-variation guarantees
must be distinguished in comparing these results.
The shifted-composition analysis of
\citet[Theorem~5.6]{altschuler2026shifted} gives a relative-entropy
guarantee that yields a \(\widetilde O(d^{1/3})\) query bound at fixed
conditioning and total-variation accuracy. Another gradient-only route,
based on rejection sampling on diffusion path space, was developed by
\citet[Theorem~3.2(ii)]{chen2026exact}; its stated query-tail bound also
gives \(\widetilde O(d^{1/3})\) in this fixed-parameter setting.

High-order integration and polynomial collocation already reduce the
dimensional dependence of Hamiltonian Monte Carlo under stronger
regularity or structural assumptions
\citep{mangoubi2017rapid,lee2018odes}. In particular, their high-order
results use assumptions beyond the two-sided Hessian bounds in
\eqref{eq:original-class}. The obstacle here is to obtain high-order
accuracy under those bounds while implementing every conditional mean
with finitely many first-order queries.

More recently, \citet[Theorem~7.1(i)]{chen2026picard} obtain
a \(\widetilde O(d^{1/6})\) expected-gradient-query bound at fixed
relative-entropy accuracy and fixed conditioning. Their theorem applies
with the supplied minimizer in our model, and Pinsker's inequality gives
the same dimension exponent at fixed total-variation accuracy.
The expected budget can also be capped at a constant multiple of its
uniform bound, using a fixed fallback on overflow, without changing that
exponent.\footnote{For example, take their relative-entropy parameter
\(\Delta=1/20\), so that the total-variation error is at most
\(1/(20\sqrt2)\). If \(B(d)\) bounds the expected number of queries,
capping the run at \(\lceil20B(d)\rceil\) adds at most \(1/20\) to
the error by Markov's inequality and a coupling of the two runs.
The resulting error is less than \(1/10\).}
Their separate high-accuracy total-variation result combines this
construction with the proximal bouncy-particle sampler
\citep{chen2026bps} and has a \(d^{1/4}\) term, while its dependence
on inverse accuracy is polylogarithmic
\citep[Theorem~1.3]{chen2026picard}.
A subsequent preprint of September 30, 2026 by
\citet[Theorem~1.1]{chen2026cloud} improves the high-accuracy
dimension dependence to \(\widetilde O(d^{1/5})\) expected gradient
queries at fixed conditioning, while retaining polylogarithmic dependence
on inverse total-variation accuracy. Its Gaussian-cloud representation
propagates unevaluated Picard iterates using likelihood corrections.

The smoothed Picard work is also a close methodological precedent:
it develops Gaussian smoothing of \(C^2\) potentials, conditional
cumulant estimates for derivatives of the smoothed potential, and
high-order Picard approximation, including
arbitrary Picard depth when exact smoothed gradients are available.
Here we prove the all-split tensor estimates needed by our construction
and implement recursively nested noisy conditional means with finitely
many queries to the original gradient. This finite implementation,
together with the subsequent removal of the auxiliary noise, gives
every fixed positive query exponent in Theorem~\ref{thm:main}.

The Gaussian conditional-sampling framework used below was developed by
\citet{lee2021proximal}; see also the functional-inequality analysis of
\citet{chen2022proximal}. Its main operation adds a Gaussian quadratic
to the potential and samples the resulting restricted conditional law.
We retain that framework and give the contraction argument used here,
while implementing its conditional operation by the recursive
construction above.

For lower bounds, \citet{chewi2023lower} developed a systematic
first-order query-complexity treatment of log-concave sampling. Their
Gaussian lower bounds use rotational invariance to reduce arbitrary
adaptive algorithms to block-Krylov information. Their results already give logarithmic dimension dependence for
well-conditioned Gaussian sampling at a suitable fixed total-variation
accuracy. Our lower-bound proof follows that strategy, with an explicit
spectrum in \([1,2]\) and a separating statistic. We give the full argument for total-variation
threshold \(1/10\), including unrestricted randomized outputs.

\subsection{Proof strategy and principal ingredients}

The auxiliary problem comes from conditioning a target sample after a
small Gaussian perturbation. Fix $h>0$. If $X\sim\pi_V$ and $G$ is an independent
standard Gaussian, the conditional density of $X$ given
$X+\sqrt h\,G=y$ is proportional to
\[
 \exp\!\left(-V(x')-\frac{|x'-y|^2}{2h}\right).
\]
Writing $x'=y+\sqrt h\,z$ gives the potential
$|z|^2/2+F_y(z)$, where $F_y(z)=V(y+\sqrt h\,z)$ and
$\Lip(\nabla F_y)\le2h$. Thus a small $h$ makes the perturbation
weak, using one original gradient query for each evaluation of
$\nabla F_y$. We first implement this conditional operation with an
added Gaussian, then remove that noise by further conditioning.

To make the conditional operation reusable, we work with a $C^2$ perturbation
$F:\R^d\to\R$, a center $x\in\R^d$, a scale $0<r\le1$,
and the near-quadratic law
\[
                 \nu(dz)\propto
                 e^{-|z|^2/2-F(x+rz)}\,dz,
\]
where the gradient of $F$ has small Lipschitz constant.
We construct two procedures: a sample from \(\nu\)
with an added Gaussian, and the mean
\(\E_\nu\nabla F(x+rZ)\) with an added Gaussian.
The noise is not merely an error allowance. It is
used to make recursive evaluation possible.

\paragraph{Gaussian transport.}
For independent \(Z\sim\nu\) and standard Gaussian
\(N\), consider
\(Y_\rho=\rho Z+\sqrt{1-\rho^2}\,N\).
The velocity of its marginal-law transport is a
conditional gradient mean, hence another instance
of the same auxiliary mean problem. This
conditional-velocity construction belongs to the
probability-flow and stochastic-interpolant viewpoint
\citep{song2021score,albergo2025interpolants}.

\paragraph{Dimension-controlled high derivatives.}
High-order approximation is useful only if each further derivative
does not introduce another power of the dimension. Derivatives of a
Gaussian conditional mean form tensors with many coordinate slots.
We bound each tensor as a matrix under every division of those slots
into two nonempty groups. These bounds let us estimate the contractions
arising along the flow by matrix composition. Repeated integration by
parts then incurs only polylogarithmic dimension losses, followed by
one $\sqrt d$ cost at the final vector output. The first step builds
on classical Appell polynomials and higher-order Poincar\'e estimates
\citep{anshelevich2004appell,adamczak2015concentration,kothari2017agnostic};
we prove the all-split cumulant and subsequent divergence estimates
needed here. All high derivatives come from Gaussian conditioning;
the original potential still needs only the $C^2$ regularity in
\eqref{eq:original-class}.

\paragraph{Stationary centering.}
To produce a noisy mean, we next couple one of the smoothed marginal laws above
to a Gaussian momentum law while preserving their product.
Its invariance uses a skew-matrix drift, as in
\citet[Section~2.1]{ma2015recipe}. Our transport-derived matrix
makes the momentum displacement convert a path integral of gradient
means into a Gaussian centered at the desired mean.
The derivative estimates control high-order quadrature of both the
probability transport and this stationary flow.

\paragraph{Noisy recursive evaluation.}
Fixed-depth polynomial and Picard computations
represent the transports and centering flow as
finite expressions in exact conditional means.
The difficulty is to implement nested mean
evaluations without knowing their exact centers.
Our recursive procedure has a pathwise
translation property: shifting a center can be
compensated by a specified shift of its Gaussian
random seeds, leaving every primitive query
unchanged. A deterministic covariance reduction
of those seeds then absorbs independent Gaussian
noise in the center. This converts a noisy-center
call into the correct known-center law.
There are two distinct depths: the dependency depth of a finite
expression in means, and the height of the routines called to evaluate
those means. We bound both before choosing the quadrature mesh.
Each routine is assigned a requested error and a bound on the size of
its perturbation. The update rules for these parameters make the routine
recursion finite and allow its query exponent to be arbitrarily small.
This seed transformation
is an exact finite-dimensional identity, not an
appeal to an approximate sampling oracle.

\paragraph{Return to the original target.}
A Gaussian proximal chain uses the near-quadratic
procedure to sample \(\pi_V\) with a small added
Gaussian. Repeated conditioning halves the
remaining variance and reduces the conditional
problem to an almost Gaussian law. A mode-centered
Gaussian completes the descent. This final
step is necessary because the requested accuracy
is for the unnoised target in total variation.

Section~\ref{sec:preliminaries} states the auxiliary
sampling guarantee and the elementary comparison
facts. Sections~\ref{sec:jets} and \ref{sec:divergence}
prove the derivative estimates.
Sections~\ref{sec:centering} and \ref{sec:numerics}
construct and discretize the two flows.
Section~\ref{sec:compiler} implements the resulting
expressions with actual gradient calls.
Section~\ref{sec:sampling} proves the sampling
upper bound, and Section~\ref{sec:lower} gives
the independent lower bound.
